% RiskPilot — ACM SIGCONF format course report (≤6 pages incl. refs & appendix)
% Compile: pdflatex main → bibtex main → pdflatex main × 2

\documentclass[sigconf,nonacm]{acmart}

\usepackage{booktabs}
\usepackage{graphicx}
\usepackage{amsmath}
\usepackage{balance}
\usepackage{xcolor}

% Tight layout for 6-page limit
\settopmatter{printacmref=false}
\setcopyright{none}
\renewcommand\footnotetextcopyrightpermission[1]{}

\title{RiskPilot: An Autonomous Evolution Engine for Financial Risk Control Rules via GNN--LLM Collaboration}

\author{Xiaoye Su}
\affiliation{%
  \institution{Hongkong University of Science and Technology(Guangzhou)}
  \city{Guangzhou}
  \country{China}}
\email{xsu315@connect.hkust-gz.edu.cn}

\author{Yifei Shi}
\affiliation{%
  \institution{Hongkong University of Science and Technology(Guangzhou)}
  \city{Guangzhou}
  \country{China}}
\email{yifei@example.com}

\begin{abstract}
Financial platforms rely on hand-crafted risk rules that lag behind evolving fraud tactics.
We present \textbf{RiskPilot}, a closed-loop system that bridges graph neural network (GNN) perception and large language model (LLM) rule synthesis.
Stage~1 trains Beta Wavelet GNN (BWGNN)~\cite{tang2022rethinking} for anomaly detection and extracts structured risk insights via embedding clustering.
Stage~2 uses retrieval-augmented generation (RAG)~\cite{lewis2020retrieval} and chain-of-thought prompting to generate executable JSON rules, with greedy threshold search and iterative refinement.
Stage~3 backtests rules on labeled graphs; Stage~4 deposits qualified rules into a vector knowledge base.
On T-Finance, Yelp, and Amazon fraud graphs, BWGNN outperforms four classical baselines by 11.9--37.4\% Macro-F1.
The full pipeline yields up to 79.8\% rule recall and 6 qualified rules on Amazon; ablation shows that \texttt{gnn\_anomaly\_score} is essential---without it, T-Finance recall drops from 69.9\% to 0.3\%.
\end{abstract}

\keywords{Graph neural networks, anomaly detection, financial risk control, large language models, rule generation}

\begin{document}
\maketitle

% ============================================================
\section{Introduction}
\label{sec:intro}

Mobile payment fraud is increasingly organized and graph-structured: attackers form dense transaction rings while mimicking normal feature profiles.
Two toolchains address complementary gaps. GNNs detect structural anomalies but output scores and embeddings, not deployable policies.
LLMs generate readable rules but lack graph perception and tend to emit arbitrary thresholds.

\textbf{RiskPilot} closes this loop in four stages: (1)~GNN perception and insight extraction; (2)~LLM rule generation with RAG; (3)~sandbox backtesting; (4)~knowledge-base accumulation with feedback.
Our contributions are:
\begin{enumerate}
  \item A \textbf{GNN-to-LLM insight bridge} that clusters anomaly embeddings and emits JSON risk patterns (Section~\ref{sec:method}).
  \item An \textbf{autonomous rule evolution loop} with threshold search, targeted LLM feedback, and early-stop stability.
  \item \textbf{Multi-dataset evaluation} on T-Finance, Yelp, and Amazon with baselines, ablations, and component analysis.
\end{enumerate}

% ============================================================
\section{Related Work}
\label{sec:related}

\textbf{Graph anomaly detection.}
Spatial GNNs (GCN~\cite{kipf2017semi}, GAT~\cite{velickovic2018graph}, GraphSAGE~\cite{hamilton2017inductive}) aggregate neighborhood signals but may miss spectral anomalies.
BWGNN~\cite{tang2022rethinking} uses Beta wavelet filters on the graph Laplacian for multi-scale spectral analysis and achieves strong results on financial graphs~\cite{liu2022bond}.

\textbf{LLM for security and RAG.}
LLMs have been applied to threat summarization and policy drafting, yet graph-structural fraud understanding and executable rule synthesis remain under-explored.
RAG~\cite{lewis2020retrieval} supplies historical rules to constrain generation quality.

% ============================================================
\section{Methodology}
\label{sec:method}

\subsection{System Overview}
RiskPilot executes \textit{GNN Perception} $\rightarrow$ \textit{LLM Strategy} $\rightarrow$ \textit{Sandbox Validation} $\rightarrow$ \textit{Rule KB}, with RAG feeding prior rules back into generation.

\subsection{Stage 1: GNN Perception}
Given graph $G{=}(V,E)$, features $X$, and normalized Laplacian $\hat{L}$, BWGNN applies $d{+}1$ Beta wavelet bases
$\psi_i(\lambda) = \frac{(\lambda/2)^i(1-\lambda/2)^{d-i}}{B(i{+}1,d{+}1-i)}$,
producing anomaly probabilities $\hat{y}$ and embeddings $Z$.

\textbf{Insight extraction.} Top-$K$ nodes ($K{=}500$, $\tau{=}0.5$) are clustered (KMeans, $C{=}5$).
For each cluster we compute degree statistics, local subgraph density, feature $z$-scores, and a natural-language \texttt{risk\_description} consumed by the LLM.

\subsection{Stage 2: LLM Strategy}
The LLM (Qwen-Plus via DashScope) receives risk insights, data-distribution statistics, and RAG-retrieved rules.
It outputs JSON rules with \texttt{conditions} over fields such as \texttt{gnn\_anomaly\_score}, \texttt{embedding\_cluster\_id}, \texttt{node\_degree\_zscore}, and \texttt{feature\_dim\_X\_zscore}.
Chain-of-thought prompts guide distribution analysis before threshold setting.
After generation, \textbf{threshold search} greedily adjusts each threshold within $\pm 30\%$ to maximize $F_1 - 0.5 \cdot \max(0, \text{FPR} - 0.05)$.
Up to 5 \textbf{closed-loop iterations} apply targeted feedback (threshold change $\leq 20\%$) with early stop after two non-improving rounds.

\subsection{Stages 3--4: Sandbox and Knowledge Base}
Rules compile to Python predicates and run on the full labeled graph with injected GNN outputs.
A rule \textit{qualifies} if $F_1 \geq 0.6$, Recall $\geq 0.5$, Precision $\geq 0.4$, and FPR $\leq 0.05$.
Qualified rules are embedded into ChromaDB for RAG in subsequent rounds.

% ============================================================
\section{Experiments}
\label{sec:exp}

\subsection{Setup}
We evaluate on three public fraud graphs (40\% train / 20\% val / 40\% test, stratified).
BWGNN trains for 100 epochs ($d{=}2$, hidden 64) on CPU.
Table~\ref{tab:datasets} summarizes datasets; all experiments use identical splits for fair comparison.

\begin{table}[t]
\caption{Dataset statistics.}
\label{tab:datasets}
\small
\begin{tabular}{lrrrr}
\toprule
Dataset & Nodes & Edges & Feat. & Anomaly \\
\midrule
T-Finance & 39{,}357 & 42.4M & 10 & 4.58\% \\
Yelp      & 45{,}954 &  8.1M & 32 & 14.53\% \\
Amazon    & 11{,}944 &  9.6M & 25 & 6.87\% \\
\bottomrule
\end{tabular}
\end{table}

\subsection{GNN Anomaly Detection}
Table~\ref{tab:gnn} reports test Macro-F1 and AUC.
BWGNN substantially exceeds Isolation Forest, LOF, One-Class SVM, and PCA+Mahalanobis on all datasets (details in Appendix~\ref{app:baselines}).

\begin{table}[t]
\caption{BWGNN vs.\ best classical baseline (Macro-F1 / AUC).}
\label{tab:gnn}
\small
\begin{tabular}{lccccc}
\toprule
Dataset & BWGNN F1 & BWGNN AUC & Best Trad. & Trad. F1 & $\Delta$F1 \\
\midrule
T-Finance & 90.4\% & 96.0\% & OCSVM  & 53.0\% & +37.4 \\
Yelp      & 70.7\% & 83.1\% & OCSVM  & 58.8\% & +11.9 \\
Amazon    & 92.4\% & 96.5\% & OCSVM  & 72.0\% & +20.3 \\
\bottomrule
\end{tabular}
\end{table}

\subsection{Risk Pattern Discovery}
On T-Finance, the insight extractor finds five patterns among 500 high-score nodes.
Table~\ref{tab:clusters} shows representative clusters; all exhibit high connectivity and near-unity anomaly scores, consistent with organized fraud rings.

\begin{table}[t]
\caption{T-Finance risk clusters (top-500 anomalies).}
\label{tab:clusters}
\small
\begin{tabular}{crrrl}
\toprule
Cl. & Nodes & Avg Deg. & Density & Pattern \\
\midrule
0 & 296 & 593  & 0.45 & Large organized ring \\
1 &  10 & 1276 & 0.71 & Ultra-dense micro-gang \\
2 &  35 & 1127 & 0.92 & High-cohesion clique \\
3 &   6 & 1318 & 1.00 & Perfect-density cell \\
4 & 153 & 924  & 0.77 & Mid-scale coordinated \\
\bottomrule
\end{tabular}
\end{table}

\subsection{End-to-End Rule Generation}
Table~\ref{tab:pipeline} summarizes the full RiskPilot pipeline.
Amazon achieves the best rule quality (79.8\% combined recall, 0.75\% FPR, 6 qualified rules).
Yelp yields no qualified rules because weaker GNN detection (70.7\% F1) limits usable \texttt{gnn\_anomaly\_score} signals.
On T-Finance, closed-loop iteration improves qualified rules from 0 (initial round) to 5 deposited rules; best-round combined recall reaches 81.5\% (Appendix~\ref{app:iteration}).

\begin{table}[t]
\caption{Full RiskPilot pipeline results.}
\label{tab:pipeline}
\small
\begin{tabular}{lccccr}
\toprule
Dataset & GNN F1 & Recall & Prec. & FPR & Qual. \\
\midrule
T-Finance & 90.4\% & 69.9\% & 76.0\% & 3.50\% & 5 \\
Amazon    & 92.4\% & 79.8\% & 88.8\% & 0.75\% & 6 \\
Yelp      & 70.7\% & 24.6\% & 39.5\% & 6.42\% & 0 \\
\bottomrule
\end{tabular}
\end{table}

\subsection{Component Analysis}
Table~\ref{tab:ablation} compares variants on T-Finance and Amazon.
\textbf{Statistical rules} (mechanical $z$-score conjunctions) achieve only 61.3\% recall and 2 qualified rules vs.\ 69.9\% and 5 for LLM rules---LLM condition selection is critical.
\textbf{Single-round generation} (no iteration) reaches 85.3\% recall but only 2 qualified rules; iteration increases qualified count by 150\% while controlling FPR.
\textbf{No-GNN-field ablation} collapses T-Finance recall to 0.3\% and Amazon recall to 49.2\%, confirming that rules are interpretable wrappers around GNN scores.

\begin{table}[t]
\caption{Component ablation (combined recall / qualified rules).}
\label{tab:ablation}
\small
\begin{tabular}{lcc|cc}
\toprule
 & \multicolumn{2}{c|}{T-Finance} & \multicolumn{2}{c}{Amazon} \\
Variant & Recall & Qual. & Recall & Qual. \\
\midrule
RiskPilot (full)     & 69.9\% & 5 & 79.8\% & 6 \\
Statistical rules    & 61.3\% & 2 & 33.1\% & 0 \\
No iteration (1 round) & 85.3\% & 2 & --- & --- \\
No GNN fields        &  0.3\% & 0 & 49.2\% & 0 \\
\bottomrule
\end{tabular}
\end{table}

% ============================================================
\section{Discussion and Conclusion}
\label{sec:disc}

\textbf{Findings.}
(1)~Graph structure is indispensable: BWGNN leads classical methods by large margins, especially on T-Finance (+37.4\% F1).
(2)~Rule quality is bounded by GNN quality---Yelp's weaker detector prevents any rule from passing sandbox gates.
(3)~The GNN-to-LLM bridge is the key enabler: without \texttt{gnn\_anomaly\_score}, rules reduce to ineffective feature thresholds.
(4)~Closed-loop iteration improves \textit{rule portfolio quality} (count and diversity), not just single-round recall.

\textbf{Limitations.}
Results depend on LLM API availability; KMeans assumes spherical clusters; temporal dynamics are not modeled; T-Social (5.7M nodes) requires mini-batch training.

\textbf{Conclusion.}
RiskPilot demonstrates that GNN perception and LLM rule synthesis can be coupled into an autonomous, self-improving risk-control pipeline.
By translating spectral graph anomalies into executable, backtested rules, it offers a practical path toward data-driven fraud policy evolution.

% ============================================================
\balance
\bibliographystyle{ACM-Reference-Format}
\bibliography{references}

% ============================================================
\appendix
\section{Appendix}
\label{app:main}

\subsection{Classical Baseline Details}
\label{app:baselines}
Table~\ref{tab:trad} lists full results on T-Finance (representative; Yelp and Amazon follow the same trend).

\begin{table}[h]
\caption{T-Finance classical baselines vs.\ BWGNN.}
\label{tab:trad}
\footnotesize
\begin{tabular}{lcccc}
\toprule
Method & Recall & Prec. & Macro-F1 & AUC \\
\midrule
Isolation Forest   & 0.1\% & 0.1\% & 47.6\% & 35.1\% \\
LOF                & 2.5\% & 2.3\% & 48.7\% & 41.6\% \\
One-Class SVM      & 32.4\% & 9.9\% & 53.0\% & 72.9\% \\
PCA + Mahalanobis  & 0.1\% & 0.1\% & 47.6\% & 41.8\% \\
\textbf{BWGNN}     & \textbf{76.0\%} & \textbf{88.2\%} & \textbf{90.4\%} & \textbf{96.0\%} \\
\bottomrule
\end{tabular}
\end{table}

\subsection{T-Finance Closed-Loop Iteration}
\label{app:iteration}
Table~\ref{tab:iter} traces iteration on T-Finance (5 rounds max, early stop enabled).

\begin{table}[h]
\caption{T-Finance iteration progress.}
\label{tab:iter}
\footnotesize
\begin{tabular}{crrrr}
\toprule
Round & Qual./Total & Recall & Prec. & FPR \\
\midrule
Initial & 0/5 & 44.0\% & 78.1\% & 0.59\% \\
1       & 1/5 & 57.7\% & 68.2\% & 1.29\% \\
2       & 2/4 & 81.6\% & 70.2\% & 1.66\% \\
3       & 0/4 & 31.4\% & 72.3\% & 0.58\% \\
4 (best)& 3/5 & 81.5\% & 71.0\% & 1.60\% \\
\midrule
Final KB & 5 rules & 69.9\% & 76.0\% & 3.50\% \\
\bottomrule
\end{tabular}
\end{table}

\subsection{Example Qualified Rule}
All qualified rules anchor on \texttt{gnn\_anomaly\_score}. Example (Amazon, Cluster~0):
\begin{verbatim}
IF gnn_anomaly_score >= 0.49
AND feature_dim_19_zscore > 3.14
THEN action = block
\end{verbatim}
Threshold search and iteration adjusted original LLM thresholds by 15--20\% to meet recall targets while keeping FPR $\leq 5\%$.

\subsection{Reproducibility}
Code, configs, and evaluation JSON are in the RiskPilot repository.
Run \texttt{python -m src.pipeline --dataset tfinance --mode full} for the full loop, or \texttt{python scripts/run\_baselines.py --all} for baselines.
LLM stage requires DashScope/OpenAI API keys; GNN and sandbox stages run without API access.

\end{document}
